 \chapter{Тодорхойгүй интеграл бодох}
\section{Рационал бутархайн интеграл}
$\dfrac{P(x)}{Q(x)}$ хэлбэрийн функцийг рационал функц гэдэг. Энд $P(x)$, $Q(x)$ нь олон гишүүнт. Хэрэв  $\deg P<\deg Q$ бол зөв рационал бутархай бусад үед засагдах рационал бутархай гэдэг. Хэрэв $\dfrac{P(x)}{Q(x)}$ засагдах рационал бутархай бол түүнийг 
$$\dfrac{P(x)}{Q(x)}=M(x)+\dfrac{R(x)}{Q(x)}$$ хэлбэрт нэг утгатай тавьж болно. Энд $M(x),\, R(x)$  олон гишүүнтүүд ба $\deg R<\deg Q$. Тиймээс зөв бутархайн интегралыг бодчихвол ямар ч рационал бутархай функцийн интегралыг бодож болно. 
\begin{theorem}[Хялбар бутархайд задлах]
$\dfrac{P(x)}{Q(x)}$ зөв рационал бутархай бол түүнийг дараах 4 хэлбэрийн бутархайнуудын нийлбэрт нэг утгатай задлаж болно. Үүнд $p^2-4q<0$ ба 
 \begin{eqnarray}
\dfrac{A}{(x-a)}&&\text{-1-р төрлийн хялбар бутархай}\nonumber\\
\dfrac{A}{(x-a)^k}&&\text{-2-р төрлийн хялбар бутархай}\nonumber\\
\dfrac{A}{(x^2+px+q)}&&\text{-3-р төрлийн хялбар бутархай}\nonumber\\
\dfrac{A}{(x^2+px+q)^k}&&\text{-4-р төрлийн хялбар бутархай}\nonumber
\end{eqnarray}
\end{theorem}
Дээрх теоремээс харахад 4 төрлийн хялбар бутархайн интегралыг бодчихвол ямар ч рационал бутархайн интеграл бодогдоно. Эхлээд эдгээр интегралуудыг хэрхэн бодох талаар үзье. 
Хялбар бутархайн интегралууд дараах байдалтай:

\[
\int \frac{dx}{x-a}=\ln|x-a|+C,
\]

\[
\int \frac{dx}{(x-a)^n}=\frac{-1}{(n-1)(x-a)^{\,n-1}}+C,
\qquad n\ge 2,
\]

\[
\int\frac{dx}{x^2+a^2}=\frac{1}{a}\arctan\frac{x}{a}+C,
\]

\[
\int\frac{dx}{(x^2+px+q)}=\frac{2}{\sqrt{4q-p^2}}
\arctan\frac{2x+p}{\sqrt{4q-p^2}}+C.
\]

Хэрэв $x^2+px+q$ квадрат нь бодит язгуургүй бол түүнийг
\[
x^2+px+q=(x+a)^2+b^2
\]
хэлбэрт оруулж арктангенс ашиглан бодно.  Иймд 4-р төрлийн хялбар бутархайн интегралийн хувьд түүнийг $t=\dfrac{2x+p}{\sqrt{4q-p^2}}$ орлуулгаар $\dfrac{A}{(t^2+1)^n}$ хэлбэрт хувиргаж болно. Иймд
$$\displaystyle I_n = \int \frac{dx}{(1+x^2)^n}$$ хэрхэн бодохыг тайлбарлая.\\

Эхлээд $n=1$ үед
\[
I_1 = \int \frac{dx}{1+x^2} = \arctan x+C
\]
байхыг харах амархан. 

Хэсэгчлэн интегралчлах замаар
\[
I_n = \frac{x}{2(n-1)(1+x^2)^{n-1}} + \frac{2n-3}{2(n-1)} I_{n-1}
\]
томьёог гаргаж болно. Иймд 4-р төрлийн бутархайн интеграл бодогдоно. Жишээ болгон эхний хэдэн интегралыг бичвэл:
\[
I_2=\frac{x}{2(1+x^2)}+\frac12 I_1
   =\frac{x}{2(1+x^2)}+\frac12\arctan x+C
\]
$I_3$–г олохын тулд рекурент томьёонд $n=3$ гэж орлуулбал
\[
I_3=\frac{x}{2(1+x^2)^2}+\frac34 I_2
\] болно.
$I_2$–ийн оронд өмнө олсоноо орлуулбал:
\[
I_3=\frac{x}{2(1+x^2)^2}
    +\frac{3}{4}\left(\frac{x}{2(1+x^2)}+\frac12\arctan x\right)+C
\]
буюу
\[
I_3=
\frac{x}{2(1+x^2)^2}
+\frac{3x}{8(1+x^2)}
+\frac{3}{8}\arctan x+C
\]
болно.
$I_4$–г ижлээр олбол
\[
I_4=
\frac{x}{2(1+x^2)^3}
+\frac{5x}{12(1+x^2)^2}
+\frac{15x}{48(1+x^2)}
+\frac{15}{48}\arctan x+C.
\]
Бүх үзсэн зүйлсээ нэгтгэн үзвэл рационал бутархайн интегралыг бодохдоо дараах алхамуудыг хийнэ. 
\begin{enumerate}
    \item Эхний алхам рационал бутархайг засаж зөв бутархай ба олон гишүүнтийн нийлбэр болгоно.
    \item Зөв бутархайг хялбар рационал бутархайнуудын нийлбэр байдлаар илэрхийлэн.
    \item Интегралаа бодно.
\end{enumerate}
$\dfrac{R(x)}{Q(X)}$ зөв бутархайн хувьд
\begin{eqnarray}
Q(x)&=&A(x-a_1)\ldots (x-a_m)(x-b_1)^{n_1}\ldots (x-b_k)^{n_k}\nonumber\\
   &=& (x^2+p_1x+q_1)\ldots (x^2+p_lx+q_l)(x^2+r_1x+s_1)^{g_1}\ldots (x^2+p_dx+q_d)^{g_d}\nonumber
\end{eqnarray} бол түүнийг хялбар бутархайд задлахад $Q(x)$ олон гишүүнтийн зэрэгтэй тэнцүү тооний үл мэдэгдэх бүхий шугман тэгшитгэлийн систем бодно.
\begin{example}
$$
\int\frac{3x^3+2x^2+x-7}{x^2+1} dx$$ хувьд

 \[
\frac{3x^3+2x^2+x-7}{x^2+1}
=
3x+2+\frac{-2x-9}{x^2+1}.
\]
болно. 
Энэ үед интеграл нь:
\[
\int (3x+2)\,dx - \int\frac{2x}{x^2+1}dx - 9\int\frac{dx}{x^2+1}=\dfrac{3}{2}x^2+2x-\ln(x^2+1)-9\arctan x+c.
\]   
\end{example} 
\begin{example}
\[
\int \frac{dx}{x^2-1}
\]
бодъёо. Зөв бутархай учир хуваарийг үржигдэхүүн болговол ${x^2-1}=(x-1)(x+1)$ тул
\[
\frac{1}{(x-1)(x+1)}=\frac{A}{x-1}+\frac{B}{x+1}
\]
 байхаар $A, B$ олдоно. Ижил хуваарь өгвөл 
 \[
\frac{A}{x-1}+\frac{B}{x+1}=\frac{A(x+1)+B(x-1)}{(x-1)(x+1)}=\frac{(A+B)x+(A-B)}{(x-1)(x+1)}=\frac{1}{(x-1)(x+1)}
\] болно. Иймд 
$$
\begin{cases}
    A+B&=0\\
    A-B&=1
\end{cases}$$ болно. Үүнийг бодвол $A=0.5, B=-0.5$ болно. 
Тэгвэл интеграл
\[
\int \frac{dx}{x^2-1}=\frac12\Bigl[\int \Bigl(\dfrac{1}{x+1}-\dfrac{1}{x-1}\Bigr)dx\Bigr]==\frac12\ln|x-1|-\frac12\ln|x+1|+C
=\frac12\ln\left|\frac{x-1}{x+1}\right|+C.
\]
 болно.
\end{example}
\begin{example}
 \[
\int \frac{dx}{x^2+4x+5}
=
\int \frac{dx}{(x+2)^2+1}
=
\arctan(x+2)+C.
\]   
\end{example}
\begin{example}
Дараах интегралыг бод.
\[
I=\int\dfrac{2x^4 - x^3 - x^2 - 4x - 2}{x^3-x^2-x+1}\, dx.
\]

\textbf{1. Зөв бутархай болгох.}

Эхлээд олон гишүүнтийг үлдэгдэлтэй  хувааж зөв бутархайд шилжүүлнэ:
\[
\frac{2x^4 - x^3 - x^2 - 4x - 2}{x^3 - x^2 - x + 1}
= 2x+1 + \frac{2x^2 - 5x - 3}{x^3 - x^2 - x + 1}.
\]

Хуваарийг үржигдэхүүн болговол
\[
x^3 - x^2 - x + 1 
= x^2(x-1) -1(x-1)
= (x^2-1)(x-1)
= (x-1)^2(x+1).
\]

Тиймээс
\[
\frac{2x^4 - x^3 - x^2 - 4x - 2}{x^3 - x^2 - x + 1}
= 2x+1 + \frac{2x^2 - 5x - 3}{(x-1)^2(x+1)}.
\]


\textbf{2. Зөв бутархайг хялбар бутархайд задлах.}

Бутархайг
\[
\frac{2x^2 - 5x - 3}{(x-1)^2(x+1)}
= \frac{A}{x-1} + \frac{B}{(x-1)^2} + \frac{C}{x+1}
\]
гэж задлана.

Ижил хуваарь оруулбал:
\[
2x^2 - 5x - 3 
= A(x-1)(x+1) + B(x+1) + C(x-1)^2
= (A+C)x^2 + (B-2C)x + (-A+B+C).
\]
болно.
Коэффициентүүдийг харьцуулбал:
\[
2x^2 - 5x - 3 = (A+C)x^2 + (B-2C)x + (-A+B+C),
\]
ингээд систем:
\[
\begin{cases}
A + C = 2,\\
B - 2C = -5,\\
-\,A + B + C = -3.
\end{cases}
\]

Энэ системээс:
\[
A = 1,\qquad B = -3,\qquad C = 1.
\]

Иймээс
\[
\frac{2x^2 - 5x - 3}{(x-1)^2(x+1)}
= \frac{1}{x-1} - \frac{3}{(x-1)^2} + \frac{1}{x+1}.
\]

Тэгвэл бүхэлд нь
\[
\dfrac{2x^4 - x^3 - x^2 - 4x - 2}{x^3 - x^2 - x + 1}
= 2x+1 + \frac{1}{x-1} - \frac{3}{(x-1)^2} + \frac{1}{x+1}.
\]

\textbf{3. Хялбар бутархай тус бүрээр интегралчлах.}

Одоо
\[
\begin{aligned}
I
&= \int (2x+1)\,dx 
+ \int \frac{1}{x-1}\,dx 
- \int \frac{3}{(x-1)^2}\,dx 
+ \int\frac{1}{x+1}\,dx.
\end{aligned}
\]

\[
\int (2x+1)\,dx = x^2 + x + C_1,
\]

\[
\int\frac{1}{x-1}\,dx = \ln|x-1| + C_2,
\]

\[
\int \frac{3}{(x-1)^2}\,dx 
= 3\int (x-1)^{-2}\,dx 
= 3\cdot\frac{(x-1)^{-1}}{-1} 
= -\frac{3}{x-1},
\]
тэгэхээр
\[
-\int \frac{3}{(x-1)^2}\,dx = +\frac{3}{x-1},
\]

\[
\int\frac{1}{x+1}\,dx = \ln|x+1| + C_3.
\]

Иймээс
\[
I
= x^2 + x + \ln|x-1| + \frac{3}{x-1} + \ln|x+1| + C.
\]

Логарифмуудыг нэгтгээд
\[
I
= x^2 + x + \frac{3}{x-1} + \ln|x^2-1| + C,
\]
учир нь $x^2-1=(x-1)(x+1)$.

\end{example}
\section{Иррационал функцийн интеграл}

Иррационал интеграл гэдэг нь интегралын тэмдэгтийн дор квадрат язгуур, ерөнхийдөө $x$-ийн бутархай зэрэг (\,$\sqrt{\phantom{x}}$, $\sqrt[n]{\phantom{x}}$\,) орсон интегралуудыг хэлнэ. Жишээлбэл:
\[
\int \sqrt{x+1}\,dx,\qquad 
\int \frac{dx}{\sqrt{x^2+4}},\qquad
\int \frac{x\,dx}{\sqrt{1-x^2}},\ \text{г.м.}
\]

Ийм интегралуудыг бодоход орлуулах арга түгээмэл ашиглана.

\subsubsection*{1. $\sqrt{ax+b}$ хэлбэрийн иррационал интеграл}

Ихэвчлэн $x$–ийг язгууртай тэнцүүлэх буюу
\[
 ax+b=t^2
\]
гэсэн орлуулга хэрэглэнэ.

\begin{example}
\[
\int \sqrt{x+1}\,dx.
\]

Орлуулая:
\[
x+1 = t^2 \quad\Rightarrow\quad x = t^2-1,\quad dx = 2t\,dt.
\]
Тэгвэл
\[
\sqrt{x+1} = \sqrt{t^2} = t \quad (t>0\ \text{мужийг авч болно}).
\]

Интеграл:
\[
\int \sqrt{x+1}\,dx
= \int t \cdot 2t\,dt
= 2\int t^2\,dt
= 2\cdot \frac{t^3}{3} + C
= \frac{2}{3}t^3 + C.
\]

Энд $t^2 = x+1$ тул $t=\sqrt{x+1}$:
\[
\int \sqrt{x+1}\,dx 
= \frac{2}{3}(x+1)^{3/2} + C.
\]
\end{example}

\subsubsection*{2. $\sqrt{x^2+a^2}$ ба $\sqrt{x^2-a^2}$: тригонометр орлуулга}

Квадрат язгуур дотор тэгш өнцөгт гурвалжин үүсгэж болохоор хэлбэр байвал тригонометр орлуулга ашиглана.

\begin{itemize}
  \item $\sqrt{x^2+a^2}$ үед $x = a\tan t$ гэж орлуулна.
  \item $\sqrt{x^2-a^2}$ үед $x = a\sec t$ (эсвэл $x=a\cosh t$ гэх мэт) гэж орлуулна.
  \item $\sqrt{a^2-x^2}$ үед $x = a\sin t$ гэж орлуулна.
\end{itemize}

\begin{example}
\[
\int \frac{dx}{\sqrt{x^2+4}}.
\]

Энд $x^2+4 = x^2+2^2$ тул $x=2\tan t$ гэж орлуулвал:
\[
x = 2\tan t \quad\Rightarrow\quad dx = 2\sec^2 t\,dt.
\]
Мөн
\[
\sqrt{x^2+4}
= \sqrt{4\tan^2 t + 4}
= 2\sqrt{\tan^2 t+1}
= 2\sec t.
\]

Интеграл:
\[
\int \frac{dx}{\sqrt{x^2+4}}
= \int \frac{2\sec^2 t\,dt}{2\sec t}
= \int \sec t\,dt.
\]

$\displaystyle \int \sec t\,dt = \ln|\sec t + \tan t| + C$ тул:
\[
\int \frac{dx}{\sqrt{x^2+4}}
= \ln|\sec t + \tan t| + C.
\]

Одоо $x$-руу буцаая. $x=2\tan t \Rightarrow \tan t = \dfrac{x}{2}$. 
Тэгвэл
\[
\sec t = \sqrt{1+\tan^2 t}
= \sqrt{1+\frac{x^2}{4}}
= \frac{\sqrt{x^2+4}}{2}.
\]

Ингээд
\[
\sec t + \tan t = \frac{\sqrt{x^2+4}}{2} + \frac{x}{2}
= \frac{\sqrt{x^2+4}+x}{2}.
\]

Тогтмол $\ln 2$-г $C$-д шингээж болдог тул:
\[
\int \frac{dx}{\sqrt{x^2+4}}
= \ln\bigl|\sqrt{x^2+4}+x\bigr| + C.
\]
\end{example}

\subsubsection*{3. Дифференциал биномын интеграл (Чебышёвын теорем)}

Рациналь $m,n,p$ тоонуудын хувьд дараах интеграл
\[
\int x^{m}(a+b x^{n})^{p}\,dx
\]
ерөнхий тохиолдолд энгийн функцээр илэрхийлэгдэхгүй. Гэхдээ зөвхөн дараах
\textbf{гурван тусгай нөхцөл} биелэх үед рациональ функцийн интегралд шилжүүлж болдог. Үүнийг Чебышёвын теорем гэнэ.

\subsection*{Тохиолдол 1. $p$ — бүхэл тоо}
Хэрэв
\[
p\in \mathbb{Z},
\]
бол 
\[
x = z^{N},
\]
орлуулга хэрэглэнэ. Энд $N$ нь $m$ болон $n$–ийн  \emph{ерөнхий хуваарь}.

\subsection*{Тохиолдол 2. $\dfrac{m+1}{n}$ — бүхэл тоо}
Хэрэв
\[
\frac{m+1}{n}\in\mathbb{Z},
\]
бол дараах орлуулгыг хэрэглэнэ:
\[
a + b x^{n} = z^{N},
\]
энд $N$ нь $p$ бутархайн \emph{хуваарь}.

\subsection*{Тохиолдол 3. $\dfrac{m+1}{n}+p$ — бүхэл тоо}
Хэрэв
\[
\frac{m+1}{n} + p \in \mathbb{Z},
\]
бол дараах орлуулга хэрэглэнэ:
\[
a x^{-n} + b = z^{N},
\]
энд мөн $N$ нь $p$ бутархайн хуваарь.

\subsection*{Тайлбар}
Хэрэв $n=1$ бол дээрх гурван нөхцөл нь дараахтай тэнцэнэ:
\begin{enumerate}
\item $p$ — бүхэл;
\item $m$ — бүхэл;
\item $m+p$ — бүхэл.
\end{enumerate}

\subsection*{Жишээ}
Дараах интегралыг бодъё:
\[
\int x^3(2+x^4)^{1/2}\,dx.
\]

Энд $m=3$, $n=4$, $p=\frac12$. Шалгах:
\[
\frac{m+1}{n} = \frac{4}{4} = 1 \in \mathbb{Z},
\]
тиймээс \emph{Тохиолдол 2} биелж байна.

Орлуулга хийвэл:
\[
2 + x^4 = t^2, \qquad t = \sqrt{2+x^4}.
\]
Мөн
\[
4x^3\,dx = 2t\,dt \quad\Rightarrow\quad x^3dx = \frac{t}{2}\,dt.
\]

Интегралыг орлуулж бодъё:
\[
\int x^3(2+x^4)^{1/2}\,dx
= \int \left( \frac{t}{2}\right)t\,dt
= \frac12\int t^2 \, dt
= \frac{t^3}{6} + C.
\]

Буцааж орлуулбал:
\[
\boxed{
\int x^3(2+x^4)^{1/2}\,dx
= \frac{(2+x^4)^{3/2}}{6} + C.
}
\]
\newpage
\section{Семинар}
\begin{enumerate}
\item Рационал бутархайн интеграл бод.
\begin{enumerate}
    
\begin{multicols}{2}
\item $\displaystyle \int \frac{3x+1}{x^2+1}\,dx$
\item $\displaystyle \int \frac{2x-3}{x^2-3x+2}\,dx$
\item $\displaystyle \int \frac{x}{(x-1)^2(x+1)}\,dx$
\item $\displaystyle \int \frac{x}{(x-1)^2(x^2+2x+2)}\,dx$
\item $\displaystyle \int \frac{2x+3}{(x-2)(x+5)}\,dx$
\item $\displaystyle \int \frac{x}{(x+1)(x+2)(x+3)}\,dx$
\item $\displaystyle \int \frac{x^{10}}{x^2+x-2}\,dx$
\item $\displaystyle \int \frac{x^3+1}{x^3-5x^2+6x}\,dx$
\item $\displaystyle \int \frac{x^2+3x-2}{(x-1)(x^2+x+1)^2}\,dx$
\item $\displaystyle \int \frac{x^3+1}{(x^2+1)^2}\,dx$
\end{multicols}
\end{enumerate}
\item Иррационал функцийн интеграл бод.
\begin{enumerate}
\begin{multicols}{2}
\item $\displaystyle \int \frac{dx}{1+\sqrt{x}} $
\item $\displaystyle \int \frac{1-\sqrt{x+1}}{1+\sqrt[3]{x+1}}\,dx$ 
\item $\displaystyle \int \sqrt{x^3+x^2}\, dx$
\item $\displaystyle \int \frac{\sqrt{x}}{(1+\sqrt[3]{x})^2}\,dx$  
\item $\displaystyle \int \sqrt[3]{3x-x^3}\, dx$
\item $\displaystyle \int \frac{x^2}{\sqrt{1+x+x^2}} \,dx$
\item $\displaystyle \int \frac{dx}{(1-x^2)\sqrt{1-x^2}} \,dx$
\end{multicols}
\end{enumerate}
\item Тригонометрийн интеграл бод. 
\begin{enumerate}
\begin{multicols}{2}
\item $\displaystyle \int \cos^5 x\, dx$
\item $\displaystyle \int \sin ^4 x\, dx$
\item $\displaystyle \int \cos^5 x\sin^5 x\, dx$
\item $\displaystyle \int \frac{\sin^3 x}{\cos ^4x}\,dx$ 
\item $\displaystyle \int \frac{dx}{\sin^3x}\,dx$
\item $\displaystyle \int \sin 3x\cos x\,dx$
\item $\displaystyle \int \cos x\cos 2x\cos 3x\,dx$
\item $\displaystyle \int \frac{dx}{2\sin x-\cos x+5}\,dx$ 
\end{multicols}
\end{enumerate}
\end{enumerate}
 \newpage
 \section{Гэрийн даалгавар} Тохиромжтой аргаар бод.
\begin{enumerate}
\begin{multicols}{2}
\item $\displaystyle \int \frac{dx}{x^6(1+x^2)}.$

\item $\displaystyle \int \frac{x^2}{(1-x^2)^3}\,dx.$

\item $\displaystyle \int \frac{dx}{1+x^4+x^8}.$

\item $\displaystyle \int \frac{dx}{\sqrt{x+\sqrt[3]{x}}}.$

% ----------- 51–55 -----------
\item $\displaystyle \int x^2\sqrt{\frac{x}{1-x}}\,dx.$

\item $\displaystyle \int \frac{x+2}{x^2\sqrt{1-x^2}}\,dx.$

\item $\displaystyle \int \sqrt{\frac{x}{1-x}}\,dx.$
\item $\displaystyle \int \frac{dx}{\sqrt[3]{x^2(1-x)}}.$

\item $\displaystyle \int \frac{dx}{x\sqrt{1+x^3+x^6}}.$

% ----------- 11–15 -----------
\item $\displaystyle \int x\sqrt{x^4-2x^2-1}\,dx.$

\item $\displaystyle \int \frac{1+\sqrt{\frac{1-x}{1+x}}}{1-\sqrt{\frac{1-x}{1+x}}}\,dx.$

\item $\displaystyle \int \frac{(1+x)\,dx}{\sqrt{1+x^2}}.$

\item $\displaystyle \int \frac{\ln(1+x+x^2)}{(1+x)^2}\,dx.$

\item $\displaystyle \int (2x+3)\arccos(2x-3)\,dx.$

% ----------- 16–20 -----------
\item $\displaystyle \int x\ln(4+x^4)\,dx.$

\item $\displaystyle \int \frac{\arcsin x}{x^2}\cdot\frac{1+x^2}{\sqrt{1-x^2}}\,dx.$

\item $\displaystyle \int \frac{x\ln\!\left(1+\sqrt{1+x^2}\,\right)}{\sqrt{1+x^2}}\,dx.$

\item $\displaystyle \int x\sqrt{x^2+1}\,\ln\!\sqrt{x^2-1}\,dx.$

\item $\displaystyle \int \frac{x}{\sqrt{1-x^2}}\ln\frac{x}{\sqrt{1-x}}\,dx.$

% ----------- 21–25 -----------
\item $\displaystyle \int \frac{dx}{(2+\sin x)^2}.$

\item $\displaystyle \int \frac{\sin 4x}{\sin^8 x+\cos^8 x}\,dx.$

\item $\displaystyle \int \frac{\sin x}{\sqrt{1+\cos x}}\,dx.$

\item $\displaystyle \int \frac{ax^2+b}{x^2+1}\arctan x\,dx.$

\item $\displaystyle \int \frac{ax^2+b}{x^2-1}\ln\!\left|\frac{x-1}{x+1}\right|dx.$

% ----------- 26–30 -----------
\item $\displaystyle \int \frac{x\ln x}{(1+x^2)^2}\,dx.$

\item $\displaystyle \int \frac{x\arctan x}{\sqrt{1+x^2}}\,dx.$

\item $\displaystyle \int \frac{\sin 2x}{\sqrt{1+\cos^4 x}}\,dx.$

\item $\displaystyle \int \frac{x^3\arccos x}{\sqrt{1-x^2}}\,dx.$

\item $\displaystyle \int \frac{x^4\arctan x}{1+x^2}\,dx.$

% ----------- 31–35 -----------
\item $\displaystyle \int \frac{x\arctan x}{(1+x^2)^3}\,dx.$

\item $\displaystyle \int \frac{x\ln\!\left(x+\sqrt{1+x^2}\right)}{(1-x^2)^2}\,dx.$

\item $\displaystyle \int \sqrt{1-x^2}\,\arcsin x\,dx.$

\item $\displaystyle \int x(1+x^2)\arctan x\,dx.$

\item $\displaystyle \int x^x(1+\ln x)\,dx.$

% ----------- 36–40 -----------
\item $\displaystyle \int \frac{\arcsin(e^x)}{e^x}\,dx.$

\item $\displaystyle \int \frac{\arctan \frac{x}{2}}{e^x(1+e^x)}\,dx.$

\item $\displaystyle \int \frac{dx}{(e^{x+1}+1)^2-(e^{x-1}+1)^2}.$

\item $\displaystyle \int \sqrt{\tanh^2 x+1}\,dx.$

\item $\displaystyle \int \frac{1+\sin x}{1+\cos x}e^x\,dx.$

% ----------- 41–45 -----------
\item $\displaystyle \int |x|\,dx.$

\item $\displaystyle \int x\,|x|\,dx.$

\item $\displaystyle \int (x+|x|)^2\,dx.$

\item $\displaystyle \int (1+|x|-|1-x|)\,dx.$

\item $\displaystyle \int (e^{-1/x})\,dx.$

% ----------- 46–50 -----------
\item $\displaystyle \int \max(1, x^3)\,dx.$

\item $\displaystyle \int \frac{x^5}{\sqrt{1+x^2}}\,dx.$
\item $\displaystyle \int \varphi(x)\,dx$, \quad энд $\varphi(x)$ нь $x$ тооноос ойрхон бүхэл тоо хүртэлх зай.

\item $\displaystyle \int \lfloor x \rfloor\, |\sin \pi x|\, dx \qquad (x>0).$

\item $\displaystyle \int f(x)\,dx$, \quad энд
\[
f(x)=
\begin{cases}
1 - x^2, & |x|\le 1,\\
1 - |x|, & |x|>1.
\end{cases}
\]

\item $\displaystyle \int f(x)\,dx$, \quad энд
\[
f(x)=
\begin{cases}
1, & x<0,\\[4pt]
x+1, & 0\le x \le 1,\\[4pt]
2x, & x>1.
\end{cases}
\]

\item $\displaystyle \int x\, f''(x)\,dx.$

% ----------- 6–10 -----------
\item $\displaystyle \int f'(2x)\,dx.$

\item $f(x)$–ийг ол, хэрэв $\; f'(x^2)=\dfrac1{x}\quad (x>0).$

\item $f(x)$–ийг ол, хэрэв $\; f'(\sin^2 x)=\cos^2 x.$

\end{multicols}
\end{enumerate}